% !TEX root = main.tex

\chapter*{前言}
\addcontentsline{toc}{chapter}{前言}

\section*{何以「仿罗什风」译阿含？}

《杂阿含经》与巴利《相应部》同为早期结集之重要见证，但求那跋陀罗译（大正第 99 册）往往直译、冗复，
偶有误读，读来多碍。鸠摩罗什所译大乘经典以删繁就简、义译圆通著称——若\textbf{同一译场精神}
（非同一历史事实）面对\textbf{阿含之教}，文体当如何？

本书即此一\textbf{推想性的文体实验}：义理以巴利平行经为主据（信），
古文主文仿罗什节奏（雅），并附今译意（达）。成书不冒充罗什译经，而供学界与读者\textbf{对照检验}。

\section*{与现有文献之关系}

\begin{itemize}
  \item \textbf{大正《杂阿含经》}：底本与经目框架，非义理之最终标准。
  \item \textbf{巴利《相应部》《增支部》等}：厘定法义之主要依据。
  \item \textbf{比较研究}（如印顺法师、阿纳罗等）：本书单经校勘与之可互参；本书侧重\textbf{全经系统化译注}与\textbf{对照阅读}。
\end{itemize}

\section*{关于经序错简与插入}

求那跋陀罗译《杂阿含经》流传久远，\textbf{经号、品目（相应）、卷次}之间时有乖违，
学界通称\textbf{错简}：底本偶见\textbf{错置的卷题}（如某经末尾出现「杂阿含经卷第…」而并非该卷起始），
或\textbf{非本集经叙}被\textbf{插入}经序之中。此类问题在 CBETA、大正本中即可见，
并非本书新造；一般读者若只按「卷名／品目」找早期法义，容易误入。

\textbf{本书的处理原则：}
\begin{itemize}
  \item \textbf{不重排}经号与五十卷框架：仍依大正本传统经号（第 1--1362 经）分卷，
    以免与馆藏编号、既有研究脱节。
  \item \textbf{底本}对照栏导出时，剥除错置卷题及非正文之序题，不混入经文。
  \item \textbf{阅读}时对已知\textbf{插入}作编辑标注：经题行附
    \textit{〔T99插入·非相应经〕}；品目于必要时\textbf{与纸面品名分离}，以免误导。
  \item 已知较著者：\textbf{第 604 经}（阿育王传／施沙，杂入天相应）、
    \textbf{第 640--641 经}（法灭授记、半阿摩勒，杂入念处相应与卷二十七之首）。
    其内容属\textbf{阿育王传说}，\textbf{无}可靠巴利专经平行；
    主文仅压缩汉本，校勘说明中标为信度较低。卷二十七因此\textbf{并非}念处相应之卷，
    而始於上述插入，继以根力相应（自第 642 经起）。
\end{itemize}

\noindent 若需依巴利、《相应部》学术序读诵，请参印顺法师、宾根海默等研究及经藏所录平行经目；
本书\textbf{不}声称已恢复全典经序，而在前言与经级标记中\textbf{如实告知}。

\section*{成书状态与局限}

\begin{itemize}
  \item 1362 经已完成首译；216 经为底本略式之重建，须与 1146 经完整译本分档阅读。
  \item 主文与今译意由大语言模型在平行约束下生成并经审读，详见「出版说明·制作方式」。
  \item 末段部分经无巴利专经平行，校勘信度多为中等——依汉本与早期教理常识，宜保守看待。
  \item 本书为\textbf{研究读本}，非终审定本；欢迎依脚注与校勘说明勘误。
\end{itemize}

\section*{致读者}

若您使用本书，建议：先读【正文】与【校勘】，再核【底本】与【平行】；
引用 ◇ 所标重建经时，请注明「据平行经重建」。
电子版对照阅读器及进一步文献说明，见「文献与制作说明」。

\vfill
\begin{flushright}
《杂阿含经·仿罗什风新译》项目组\\
2026
\end{flushright}

\clearpage
